\documentstyle[%


righttag%


,11pt%


,amssymb%


,russian%               remove in Eng.


,izvru%                 Set izven% in Eng.


]{amsart}





% {,}





\newcommand{\ctg}{\operatorname{ctg}}


\newcommand{\const}{\operatorname{const}}


\newcommand{\tts}[1]{}





\theoremstyle{definition}


\theorembodyfont{\rm}


\newtheorem{alg}{\hskip\parindent Алгоритм}


\renewcommand{\baselinestretch}{1.06}








%\hoffset=-15mm%                  For Eng.


\setcounter{page}{62}


\god{2004}


\nomer{8~(507)} %                        Remove in Eng.


%\nomer{Vol.\,48, No.\,8}%               Set in Eng.


%\str{\pageref{begin}--\pageref{end}}%  Set in Eng.


\udk{517.968:519.642}





\author{\it В.С.\,Сизиков, А.В.\,Смирнов, Б.А.\,Федоров}


% NB:   ^^ remove in Eng.


\title{Численное решение


сингулярного интегрального уравнения Абеля


обобщенным методом квадратур}


\thanks{Работа выполнена при финансовой поддержке Министерства 


образования Российской Федерации, проект \No\,E00-3.2-387.}





\date{}


%\pagestyle{myheadings}%         Set in Eng.


\pagestyle{plain} %              Remove in Eng.


\begin{document}


%\thispagestyle{plain}%          Set in Eng.


\maketitle


\label{begin}








\section{Введение}





Вопросу численного решения сингулярных интегральных уравнений (СИУ)


посвящено большое число публикаций (\cite{iv1}--\cite{ra} и др.).


Исследованы следующие СИУ:





\noindent{}a) одномерные СИУ I-го рода с ядром Коши


\cite{bl}, \cite{gab}, \cite{mus},


например, (\cite{bl}, сс.\,8, 73),


\begin{equation}


\label{eq1}


\int_a^b \frac{y(s)\,ds}{x-s}=f(x), \quad a\le x\le b,


\end{equation}


где $y(s)$ --- искомая функция;





\noindent{}b) СИУ I-го рода с ядром Гильберта \cite{bl}, \cite{gab}


или с обобщенным ядром Гильберта \cite{gab};





\noindent{}c) СИУ II-го рода с ядрами Коши или Гильберта


\cite{iv1}, \cite{bk}, \cite{mus}, например, \cite{bk},


\begin{equation}


\label{eq2}


A(x)y(x)+\frac{B(x)}{2\pi}\int_0^{2\pi}\ctg\frac{x-s}2y(s)ds+


\frac1{2\pi}\int_0^{2\pi}h(x,s)y(s)ds=f(x),\ \ 0\le x<2\pi,


\end{equation}


где $A(x)$, $B(x)$ --- известные непрерывные функции, а $h(x,s)$


--- известная $2\pi$-периодическая функция;


если $A(x)=0$ при некоторых, но не всех значениях $x$, то


уравнение (\ref{eq2}) будет уравнением III-го рода (ср. \cite{vs}, с.\,140);





\noindent{}d) СИУ I-го или II-го рода с логарифмическими и другими (слабо


сингулярными) ядрами, например, (\cite{gab}, с.\,6),


\begin{equation}


\label{eq3}


-\frac1{2\pi}\int_0^{2\pi}\ln\Big|\sin\frac{x-s}2\Big|\,y(s)\,ds+


\frac1{2\pi}\int_0^{2\pi}h(x,s)\,y(s)\,ds=f(x), \quad 0\le x<2\pi.


\end{equation}





При этом сингулярные интегралы понимаются либо как несобственные,


либо в смысле главного значения по Коши--Лебегу.





Рассмотрены также двумерные СИУ \cite{iv1}, \cite{bk}, нелинейные


СИУ \cite{ra} и др.





Что касается численных методов решения СИУ, то при их построении


использованы следующие основные алгоритмы. 





\begin{alg}[{\series{m}\shape{n}\selectfont сдвиг сеток узлов}]


В работах (\cite{bl}, с.\,8; \cite{gab}, с.\,27; \cite{bk})


вводятся дискретные сетки узлов по переменным $s$ и $x$,


например, для уравнения~(\ref{eq1}):


$s_j=a+jh$, $x_i=s_i+\Delta$, $j,i=0,1,\dotsc,n$, $s_n=b$,


где $h=(b-a)/n$ --- шаг дискретизации, $\Delta$ --- сдвиг между сетками,


полагаемый равным $\Delta=h/2$ (\cite{bl}, с.\,8) или $\Delta\in[0,h/2]$


(\cite{gab}, с.\,27) или $\Delta\in(0,h/2)$ \cite{bk}.


Введение сдвига $\Delta$ позволяет устранить особенности при


использовании численных методов, например, метода квадратур.


\end{alg}





\begin{alg}[{\series{m}\shape{n}\selectfont метод квадратур}]


В работе (\cite{bl}, с.\,8) при построении метода дискретных вихрей


(одного из вариантов метода квадратур) интеграл в (\ref{eq1})


расписывается по формуле левых прямоугольников, причем полагается,


что на каждом промежутке $[s_j,s_{j+1})$ вся подинтегральная функция


$\frac{y(s)}{x-s}=\const=\frac{y(s_j)}{x-s_j}$


при некотором фиксированном $x=x_i$.


В результате при использовании дискретных сеток по $s$ и $x$ и сдвига


$\Delta$ между ними получается система линейных алгебраических уравнений


(СЛАУ) относительно значений $y_j\equiv y(s_j)$, матрица которой имеет


преобладающую диагональ.


В работе $\cite{bk}$ при решении уравнения (\ref{eq2}) интеграл


$\int\limits_0^{2\pi}h(x,s)\,y(s)\,ds$ также расписывается по формуле левых


прямоугольников и полагается


$h(x,s)\,y(s)=\const=h(x,s_j)\,y(s_j)$ при $s\in[s_j,s_{j+1})$,


интеграл же $\int\limits_0^{2\pi}\ctg\frac{x-s}2y(s)ds$ расписывается


по обобщенной формуле левых прямоугольников (ср. \cite{ksh}, с.\,57),


согласно которой лишь $y(s)=\const=y(s_j)$ при $s\in[s_j,s_{j+1})$,


а интеграл $\int\limits_{s_j}^{s_{j+1}}\ctg\frac{x_i-s}2ds$ вычисляется


аналитически.


При этом, чтобы он не обращался в бесконечность ни при какой


комбинации $i$ и $j$, использован сдвиг сеток $\Delta\in(0,h/2)$.


\end{alg}





\begin{alg}[{\series{m}\shape{n}\selectfont аппроксимация решения}]


В работе (\cite{gab}, сс.\,15, 27), а также (\cite{iv1}, сс.\,149, 159)


искомое решение $y(s)$ аппроксимируется алгебраическим или


тригонометрическим многочленом $y_m(s)$ порядка $m$ или


полиномиальным сплайном.


При этом, например, в (\ref{eq3}) интеграл


$\int\limits_0^{2\pi}h(x,s)y_m(s)ds$ заменяется конечной суммой по


формуле левых прямоугольников, а интеграл


$\int\limits_0^{2\pi}\ln|\sin\frac{x-s}2|y_m(s)ds$


вычисляется аналитически.


В результате коэффициенты аппроксимирующего многочлена $y_m(s)$


определяются из условия минимума невязки, что приводит к тому или


иному проекционному методу (методу квадратур, наименьших квадратов,


Гал\"еркина, коллокации, сплайнов и т.\,д.) и к решению соответствующей


СЛАУ. 


\end{alg}





Во всех перечисленных методах имеет место {\it саморегуляризация},


причем роль параметра регуляризации играет $\Delta$.


При этом чем больше $\Delta\in(0,h/2]$, тем устойчивее решение $y$,


но меньше разрешающая способность численного метода, и чем меньше


$\Delta$, тем менее устойчиво решение, но выше разрешение. 





В данной работе предлагается иной вариант численного решения


некоторых СИУ. А именно, сетки узлов по $x$ и по $s$ полагаются


совпадающими (т.\,е. $\Delta=0$), а особенности устраняются за счет


использования обобщенной квадратурной формулы (ср. \cite{ksh}, с.\,57).


Такая методика может быть применена не для всех СИУ.


Она не применима для СИУ с ядрами Коши или Гильберта, но применима


для некоторых СИУ с логарифмическими и другими (слабо сингулярными)


ядрами (\cite{gab}, сс.\,6, 232).


В данной работе рассматриваются лишь СИУ I-го рода с ядром вида


$K(x,s)=s/\sqrt{s^2-x^2}$ и с переменным нижним пределом


(уравнение Абеля) и метод квадратур (обобщенный метод квадратур)


его решения.








\section{Метод численного решения СИУ\protect\\ на основе 


обобщенной квадратурной формулы}





Рассмотрим сингулярное интегральное уравнение Абеля (уравнение I-го рода


с переменным нижним пределом) в форме\footnote{Уравнение Абеля находит


широкое применение в прикладных задачах диагностики плазмы, астрофизики,


оптической и рентгеновской дифракции и т.\,д. (\cite{vs}, с.\,109;


\cite{siz}, с.\,97--98). В астрофизике оно часто называется уравнением


Цейпеля (\cite{vs}, с.\,109).}


\begin{equation}


\label{eq4}


\int_x^R\frac{s}{\sqrt{s^2-x^2}}y(s)ds=f(x),\quad 0\le x\le R,


\end{equation}


где $y(s)$ -- искомая функция и, в частности, $R=\infty$.


Заметим, что в традиционной форме уравнение Абеля записывается как


(\cite{vs}, с.\,107; \cite{k}, с.\,19)


\begin{equation}


\label{eq5}


\int_0^x\frac{y(s)}{\sqrt{x-s}}ds=2f(x),\quad x>0,


\end{equation}


причем между (\ref{eq4}) и (\ref{eq5}) имеет место


взаимно однозначное соответствие.





Уравнение (\ref{eq4}) является сингулярным уравнением со слабой


особенностью.


Данное уравнение имеет аналитическое решение (\cite{siz}, с.\,98)


\begin{equation}


\label{eq7}


y(s)=-\frac2{\pi}\int_s^R\frac{f'(x)}{\sqrt{x^2-s^2}}dx,


\quad 0\le s\le R.


\end{equation}





Однако решение (\ref{eq7}) содержит производную $f'(x)$ от


обычно экспериментальной, а значит, зашумленной функции $f(x)$,


а задача численного дифференцирования зашумленной функции


является некорректной (\cite{ta}, с.\,18--19).


Кроме того, интеграл в правой части (\ref{eq7}) является несобственным.


Поэтому задача вычисления решения (\ref{eq7}) не является тривиальной.


В данной работе предлагаются два численных метода получения решения


уравнения Абеля. 


\medskip





{\bf 1-й метод} (метод квадратур численного решения


уравнения (\ref{eq4}) на основе обобщенной формулы левых


прямоугольников --- обобщенный метод квадратур).


Введем равномерные совпадающие сетки узлов по $x$ и $s$


\begin{equation}


\label{eq8}


x_i=ih, \quad s_j=x_j=jh, \quad i,j=0,1,\dotsc,n, \quad x_n=s_n=R,


\end{equation}


где $h=\Delta x=\Delta s=R/n=\const$ --- шаг дискретизации


($R$ считаем конечным).


На некотором промежутке $[s_j,s_{j+1})$, $j\in[0,n-1]$, полагаем


\begin{equation}


\label{eq9}


y(s)=y(s_j)=\const.


\end{equation}





Легко проверяется





\begin{lem}


При условии $(\ref{eq9})$ имеет место равенство


\begin{equation}


\label{eq10}


\int_{s_j}^{s_{j+1}}\frac{s}{\sqrt{s^2-x^2}}\,y(s)\,ds=


\big(\sqrt{s_{j+1}^2-x^2}-\sqrt{s_j^2-x^2}\,\big)y_j,


\quad j\in[0,n-1],


\end{equation}


где $y_j\equiv y(s_j)$, $x\le s_j<s_{j+1}\le R$.


\end{lem}








\begin{defn}


Назовем формулу (\ref{eq10}) {\it обобщенной


формулой левых прямоугольников} (ср. \cite{ksh}, сс.\,57, 74)


для сингулярности $\frac{s}{\sqrt{s^2-x^2}}$, а множители


$\big(\sqrt{s_{j+1}^2-x^2}-\sqrt{s_j^2-x^2}\,\big)$ ---


{\it квадратурными коэффициентами} этой формулы.


\end{defn}





\begin{thm}


Численное решение уравнения $(\ref{eq4})$ по $1$-му методу


выражается следующими рекуррентными


формулами\/${:}$


\begin{equation}


\label{eq11}


\begin{split}


&y_n=y_{n-1}=\frac{f_{n-1}}{p_{n-1,n-1}}=


\dfrac{f_{n-1}}{\sqrt{R^2-(R-h)^2}},\\


&y_i=\frac{f_i-\sum\limits_{j=i+1}^{n-1}p_{ij}\,y_j}{p_{ii}},


\quad i=n-2,n-3,\dotsc,0,


\end{split}


\end{equation}


где $f_i\equiv f(x_i)$,


а


\begin{equation}


\label{eq12}


p_{ij}=\sqrt{s_{j+1}^2-x_i^2}-\sqrt{s_j^2-x_i^2}\,.


\end{equation}


\end{thm}





\begin{pf}


Интеграл в левой части (\ref{eq4}) равен сумме


интегралов (\ref{eq10}) по отдельным промежуткам $[s_j,s_{j+1})$:


\begin{equation}


\label{eq13}


\int_{x_i}^R\frac{s}{\sqrt{s^2-x_i^2}}\,y(s)\,ds=


\sum_{j=i}^{n-1}\big(\sqrt{s_{j+1}^2-x_i^2}-


\sqrt{s_j^2-x_i^2}\,\big)y_j=f_i,


\quad i=0,1,\dotsc,n-1.


\end{equation}


Соотношения (\ref{eq13}) образуют СЛАУ относительно значений $y_j$,


$j=0,1,\dotsc,n-1$.


Поскольку матрица данной СЛАУ является верхней треугольной,


то ее решение можно найти рекуррентно.


Из (\ref{eq13}), изменяя $i$ в направлении убывания, т.\,е.


полагая $i=n-1,n-2,\dotsc,0$, получим последовательно решения для


$y_{n-1},y_{n-2},\dotsc,y_0$.


Что касается значения $y_n\equiv y(R)$, то вводим дополнительное


равенство $y_n=y_{n-1}$. В результате получим формулы (\ref{eq11}),


(\ref{eq12}).


\end{pf}








{\bf 2-й метод} (метод вычисления интеграла в выражении


(\ref{eq7}) по квадратурной формуле, учитывающей его сингулярность).


Как и в 1-м методе, вводим сетки узлов (\ref{eq8}).


Далее полагаем, что на промежутке $[x_i,x_{i+1})$, $i\in[0,n-1]$, функция


\begin{equation}


\label{eq14}


f'(x)=f'(x_i)=\const.


\end{equation}


Тогда справедлива 





\begin{lem}


При условии $(\ref{eq14})$ имеет место равенство


\begin{equation}


\label{eq15}


\int_{x_i}^{x_{i+1}}\frac{f'(x)}{\sqrt{x^2-s^2}}dx=


\ln\frac{x_{i+1}+\sqrt{x_{i+1}^2-s^2}}{x_i+\sqrt{x_i^2-s^2}}


f'_i, \quad i\in[0,n-1],


\end{equation}


где $f'_i\equiv f'(x_i)$, $s\le x_i<x_{i+1}\le R$.


\end{lem}








\begin{defn}


Назовем формулу (\ref{eq15})


{\it обобщенной формулой левых прямоугольников}


для сингулярности $\frac1{\sqrt{x^2-s^2}}$, а множители


$\ln\frac{x_{i+1}+\sqrt{x_{i+1}^2-s^2}}{x_i+\sqrt{x_i^2-s^2}}$ ---


{\it квадратурными коэффициентами} этой формулы. 


\end{defn}





\begin{thm}


Пусть строится численно решение $(\ref{eq7})$,


причем интеграл в правой части $(\ref{eq7})$ вычисляется с использованием


обобщенной квадратурной формулы левых прямоугольников $(\ref{eq15})$


на сетках узлов $(\ref{eq8})$.


В этом случае решение $(\ref{eq7})$ в численном виде по $2$-му методу равно


\begin{equation}


\label{eq16}


\begin{split}


&y_0=\frac{f_0-h\sum\limits_{j=1}^{n-1}y_j}{h}, \\


&y_j=-\frac2{\pi}\sum\limits_{i=j}^{n-1}q_{ij}f'_i,


\quad j=1,2,\dotsc,n-1, \\


&y_n=y_{n-1},


\end{split}


\end{equation}


где


$$


%%\begin{equation} \label{eq17}


q_{ij}=\ln\frac{x_{i+1}+\sqrt{x_{i+1}^2-s_j^2}}{x_i+\sqrt{x_i^2-s_j^2}}.


%%\end{equation}


$$


\end{thm}





\begin{pf}


Интеграл в правой части (\ref{eq7}) равен сумме


интегралов (\ref{eq15}) по отдельным промежуткам $[x_i,x_{i+1})$:


$$


\int_{s_j}^{R}\frac{f'(x)}{\sqrt{x^2-s_j^2}}dx=


\sum_{i=j}^{n-1}\ln\frac{x_{i+1}+\sqrt{x_{i+1}^2-s_j^2}}


{x_i+\sqrt{x_i^2-s_j^2}}f'_i, \quad j=1,2,\dotsc,n-1.


$$


Отсюда получаем выражение (\ref{eq16}) для $y_j$, $j=1,2,\dotsc,n-1$.


Что касается значения $y_n$, то вводим дополнительное равенство


$y_n=y_{n-1}$.


Для получения $y_0$ формула (\ref{eq15})


не может быть использована, т.\,к. дает расходимость при $x_0=s_0=0$.


Поэтому для вычисления $y_0$ воспользуемся формулой (\ref{eq11})


1-го метода при $i=0$.


В результате получим формулы (\ref{eq16}).


\end{pf}








\section{Оценки погрешностей 1-го метода}





Выведем формулы, дающие оценки погрешностей решения уравнения (\ref{eq4})


1-м методом (ср. \cite{vs}, сс.\,42, 122). 





\begin{lem}


Интеграл


\begin{equation}


\label{eq18}


\int_{s_j}^{s_{j+1}}\frac{s}{\sqrt{s^2-x_i^2}}y(s)ds, \quad


x_i\le s_j<s_{j+1}\le R, \quad j=i,\dotsc,n-1, \quad i=0,\dotsc,n-1,


\end{equation}


при использовании обобщенной формулы левых прямоугольников $(\ref{eq10})$


и учете квадратурной погрешности может быть записан как


\begin{equation}


\label{eq19}


\int_{s_j}^{s_{j+1}}\frac{s}{\sqrt{s^2-x_i^2}}y(s)ds=


y_j\,p_{ij}+\Delta R_{ij},


\end{equation}


где $p_{ij}$ --- квадратурный коэффициент согласно $(\ref{eq12})$, а


$\Delta R_{ij}$ --- квадратурная погрешность на промежутке


$[s_j,s_{j+1})$, приближенно равная


\begin{equation}


\label{eq20}


\Delta R_{ij}=\frac{y'(\zeta_j)}2\left[\big(s_{j+1}-2s_j\big)


\sqrt{s_{j+1}^2-x_i^2}+s_j\sqrt{s_j^2-x_i^2}+


x_i^2\ln\frac{s_{j+1}+\sqrt{s_{j+1}^2-x_i^2}}


{s_j+\sqrt{s_j^2-x_i^2}}\right],


\end{equation}


причем $\zeta_j\in[s_j,s_{j+1})$.


\end{lem}





\begin{pf}


Для вычисления интеграла (\ref{eq18})


представим функцию $y(s)$ на каждом промежутке интерполяционным


полиномом Лагранжа нулевой степени $\widetilde y(s)=y_j$,


$s\in[s_j,s_{j+1})$, $j\in[0,n-1]$.


Погрешность такой интерполяции равна (\cite{bs}, с.\,504)


$\Delta y_j(s)\equiv y(s)-y_j=y'(\xi)\,(s-s_j)$,


где $\xi=\xi_j(s)\in[s_j,s_{j+1})$.


Тогда $y(s)=y_j+y'(\xi)\,(s-s_j)$ и получим (\ref{eq19}), где


$$


\Delta R_{ij}=\int_{s_j}^{s_{j+1}}\frac{s(s-s_j)}


{\sqrt{s^2-x_i^2}}y'(\xi)ds


$$


или приближенно


\begin{equation}


\label{eq21}


\Delta R_{ij}=y'(\zeta_j)\int_{s_j}^{s_{j+1}}\frac{s(s-s_j)}


{\sqrt{s^2-x_i^2}}ds,


\quad \zeta_j\in[s_j,s_{j+1}].


\end{equation}





Интеграл в (\ref{eq21}) вычисляется аналитически,


в результате получим оценку (\ref{eq20}).


\end{pf}





\begin{note}


Производную $y'(\zeta)$ можно вычислить


по приближенной формуле


\begin{equation}


\label{eq22}


y'(\zeta_j)\equiv y'_j=\begin{cases}


\frac{y_{j+1}-y_j}{h}, & j=0; \\[2mm]


\frac{y_{j+1}-y_{j-1}}{2h}, & j\in[1,n-1].


\end{cases}


\end{equation}


\end{note}





\begin{thm}


Численное решение уравнения $(\ref{eq4})$


на сетках узлов $(\ref{eq8})$ методом квадратур на основе обобщенной


формулы левых прямоугольников $(\ref{eq10})$


выражается формулами $(\ref{eq11})$,


а погрешности --- следующими рекуррентными формулами\/${:}$


\begin{equation}


\label{eq23}


\begin{split}


&\Delta y_n=\Delta y_{n-1}=\frac{\Delta R_{n-1,n-1}}{p_{n-1,n-1}}, \\


&\Delta y_i=\frac{R_i-\sum\limits_{j=i+1}^{n-1}p_{ij}\Delta y_j}


{p_{ii}}, \quad i=n-2,n-3,\dotsc,0,


\end{split}


\end{equation}


где


\begin{equation}


\label{eq24}


R_i=\sum_{j=i}^{n-1}\Delta R_{ij},


\end{equation}


причем $p_{ij}$, $\Delta R_{ij}$ и $y'(\zeta_j)$ выражаются формулами


$(\ref{eq12})$, $(\ref{eq20})$ и $(\ref{eq22})$ соответственно.


\end{thm}





\begin{pf}


Запишем интеграл в левой части (\ref{eq4})


в виде суммы интегралов по промежуткам


\begin{multline*}


\int_{x_i}^R\frac{s}{\sqrt{s^2-x_i^2}}y(s)ds=


\sum_{j=i}^{n-1}\int_{s_j}^{s_{j+1}}\frac{s}{\sqrt{s^2-x_i^2}}


\big[y_j+\Delta y_j(s)\big]ds=\\


%


=\sum_{j=i}^{n-1}\int_{s_j}^{s_{j+1}}\frac{s}{\sqrt{s^2-x_i^2}}


\big[y_j+y'(\xi)(s-s_j)\big]ds, \quad i=0,1,\dotsc,n-1,


\end{multline*}


откуда


\begin{equation}


\label{eq25}


\sum_{j=i}^{n-1}\int_{s_j}^{s_{j+1}}\frac{s}{\sqrt{s^2-x_i^2}}


\Delta y_j(s)ds=\sum_{j=i}^{n-1}\int_{s_j}^{s_{j+1}}


\frac{s(s-s_j)}{\sqrt{s^2-x_i^2}}y'(\xi)ds,


\quad i=0,1,\dotsc,n-1.


\end{equation}





Для вычисления интеграла в левой части (\ref{eq25}) используем


формулу левых прямоугольников, т.\,е. положим


$\Delta y_j(s)=\Delta y(s_j)\equiv\Delta y_j=\const$, $s\in[s_j,s_{j+1})$,


и вычислим интеграл аналогично (\ref{eq10}).


Интеграл в правой части (\ref{eq25}) равен


$\Delta R_{ij}$ согласно (\ref{eq21}). Тогда


\begin{equation}


\label{eq26}


\sum_{j=i}^{n-1}p_{ij}\Delta y_j=R_i, \quad i=0,1,\dotsc,n-1,


\end{equation}


где через $R_i$ обозначена сумма (\ref{eq24}).


Запись (\ref{eq26}) есть СЛАУ относительно $\Delta y_j$, $j=0,1,\dotsc,n-1$.


При этом полагается, что предварительно вычислены значения


$y_i$, $i=0,1,\dotsc,n$, в виде (\ref{eq11}), а затем значения


$R_i$, $i=0,1,\dotsc,n-1$, согласно (\ref{eq24}), (\ref{eq20}), (\ref{eq22}).


Решая СЛАУ (\ref{eq26}) с верхней треугольной матрицей, получим


в рекуррентной форме решение (\ref{eq23}), добавив при этом условие


$\Delta y_n=\Delta y_{n-1}$.


\end{pf}





\begin{note}


Формулы (\ref{eq23}) дают значения погрешностей решения


$\Delta y_i$ с учетом знака (ср. \cite{vs}, сс.\,42, 122),


а не по абсолютной величине вида $|\Delta y_i|\approx\dotsb$


или верхних оценок $|\Delta y_i|\le\dotsb$


или верхних оценок по норме $\|\Delta y\|\le\dotsb$


или асимптотических оценок скорости сходимости к точному решению


$\|\Delta y\|=O(\dotsb)$, как в большинстве работ (\cite{gab}, с.\,18--46


и др.).


Поэтому после получения решения $y_i$, $i=0,1,\dotsc,n$, в виде


(\ref{eq11}) и вычисления погрешностей решения $\Delta y_i$, $i=0,1,\dotsc,n$,


согласно (\ref{eq23}) можно получить уточненное решение


\begin{equation}


\label{eq27}


\widehat y_i=y_i+\Delta y_i, \quad i=0,1,\dotsc,n.


\end{equation}


\end{note}








\section{Численные примеры}





Изложенную методику численного решения СИУ обобщенным методом квадратур


(1-м методом) проиллюстрируем следующими примерами. 





\begin{exam}


Решается уравнение Абеля (\ref{eq4});


точное решение $y(s)=R-s$, $0\le s\le R$, $R=1$,


правая часть


$f(x)=\frac{R}2\sqrt{R^2-x^2}-\frac{x^2}2\ln\frac{R+\sqrt{R^2-x^2}}{x}$,


$f(0)=R^2/2$.


Сетки узлов по $s$ и по $x$: $s_j=jh$, $x_j=s_j$, $j=0,1,\dotsc,n$,


где $n=20$ --- число шагов,


$h=\Delta s=\Delta x=R/n=0{,}05$ --- шаг дискретизации. 





На рис.\,1 приведены: точное решение $y(s)$, правая часть $f(x)$,


приближенное решение $y_i$ согласно (\ref{eq11}),


погрешности решения $\Delta y_i$ в виде (\ref{eq23}) и


уточненное решение $\widehat y_i$, $i=0,1,\dotsc,n$, согласно (\ref{eq27}). 





\begin{center}


\unitlength=1mm


\begin{picture}(170,74) %% width, heigth, added 5 mm for signature


\put(32,74){\special {em:graph ssf1.bmp}} %% left X, upper Y, -, /2, +toX


\put(78,0){\mbox Рис.\,1}


\end{picture}


\end{center}


\begin{center}


--------- \quad --- точное решение $y(s)$,\quad


-- -- -- -- \quad --- правая часть $f(x)$,





$\bullet$ --- приближенное решение $y_i$,


$+$ --- погрешность решения $\Delta y_i$,


$\square$ --- уточненное решение $\widehat y_i$.


\end{center}


\end{exam}





\begin{exam}


Решается ``усеченное'' уравнение Абеля


(из задачи дифракции рентгеновских лучей)


$$


\int_q^b\frac{s}{\sqrt{s^2-q^2}}J(s)ds=f(q), \quad a\le q\le b,


$$


где $a=1$, $b=71$. Точное решение $J(s)=1/s^3$, $a\le s\le b$,


правая часть $f(q)=\sqrt{b^2-q^2}/(bq^2)$.


Сетки узлов по $s$ и по $q$: $s_j=a+jh$, $q_j=s_j$, $j=0,1,\dotsc,n$,


где $n=70$ -- число шагов, $h=\Delta s=\Delta q=(b-a)/n=1$ --- шаг


дискретизации. 





Характерная особенность примера 2 состоит в том, что диапазон изменения


как $J(s)$, так и $f(q)$ составляет несколько порядков. 





На рис.\,2 приведены: точное решение $J(s)$, функция $2f(q)$,


приближенное решение $J_i$ согласно (\ref{eq11}),


погрешности решения $\Delta J_i$ в форме (\ref{eq23}) и


уточненное решение $\widehat J_i$, $i=0,1,\dotsc,n$,


согласно (\ref{eq27}). 


\begin{center}


\unitlength=1mm


\begin{picture}(170,111) %% width, heigth, added 5 mm for signature


\put(16,111){\special {em:graph ssf2.bmp}} %% left X, upper Y, -, /2, +toX


\put(78,0){\mbox Рис.\,2}


\end{picture}


\end{center}


\begin{center}


--------- \quad --- точное решение $J(s)$,\quad


-- -- -- -- \quad --- функция $2f(q)$,





$\bullet$ --- приближенное решение $J_i$,


$+$ --- погрешность решения $\Delta J_i$,


$\square$ --- уточненное решение $\widehat J_i$.


\end{center}


\end{exam}





\begin{thebibliography}{99}





\bibitem{iv1}


Иванов В.В. {\it Теория приближенных методов и ее применение к


численному решению сингулярных интегральных уравнений.} --


Киев: Наук. думка, 1968. -- 287~с.





\bibitem{bl}


Белоцерковский С.М., Лифанов И.К. {\it Численные методы в сингулярных


интегральных уравнениях.} -- М.: Наука, 1985. -- 256~с.





\bibitem{gab}


Габдулхаев Б.Г. {\it Прямые методы решения сингулярных интегральных


уравнений первого рода.} -- Казань: Изд-во Казанск. ун-та, 1994. -- 288~с.





\bibitem{bk}


Бойков И.В., Кудряшова Н.Ю. {\it Приближенные методы решения


сингулярных интегральных уравнений в исключительных случаях} //


Дифференц. уравн. -- 2000. -- Т.\,36. -- \No\,9. -- С.\,1230--1237.





\bibitem{sam}


Самойлова Э.Н. {\it Сплайновые приближения решения сингулярного


интегродифференциального уравнения} // Изв. вузов. Математика. --


2001. -- \No\,11. -- С.\,35--45.





\bibitem{ra}


Рахимов И.К. {\it Проекционные методы решения сингулярного


интегрального уравнения Теодорсена} // Изв. вузов. Математика. --


2002. -- \No\,9. -- С.\,67--70.





\bibitem{mus}


Мусхелишвили Н.И. {\it Сингулярные интегральные уравнения.} --


3-е изд-е. -- М.: Наука, 1968. -- 512~с.





\bibitem{vs}


Верлань А.Ф., Сизиков В.С. {\it Интегральные уравнения: методы,


алгоритмы, программы.} -- Киев: Наук. думка, 1986. -- 544~с.





\bibitem{ksh}


Крылов В.И., Шульгина Л.Т. {\it Справочная книга по численному


интегрированию.} -- М.: Наука, 1966. -- 372~с.





\bibitem{siz}


Сизиков В.С. {\it Математические методы обработки результатов


измерений.} -- СПб.: Политехника, 2001. -- 240~c.





\bibitem{k}


Краснов М.Л. {\it Интегральные уравнения. Введение в теорию.} --


М.: Наука, 1975. -- 304~c.





\bibitem{ta}


Тихонов А.Н., Арсенин В.Я. {\it Методы решения некорректных задач.} --


3-е изд-е. -- М.: Наука, 1986. -- 288~c.





\bibitem{bs}


Бронштейн И.Н., Семендяев К.А. {\it Справочник по математике для


инженеров и учащихся втузов.} -- 13-е изд-е. -- М.: Наука, 1986. -- 544~c.





\end{thebibliography}





\vskip 2cm





\post{\it Санкт-Петербургский государственный}{\it Поступила}


\post{\it университет информационных технологий,}{27.02.2003}


\post{\it механики и оптики}{}





\label{end}


\end{document}





